\section{Experiments}

\subsection{Dataset and split construction}

All experiments use the \physiomio\ processing path implemented for this project~\citep{physiomioData2026,psrrepo2026}. The dataset construction step extracts contiguous movement segments from raw parquet recordings, maps gesture labels to five-bit finger targets, preprocesses each segment, and stores processed tensors as PyTorch split files.

Patients are shuffled with a fixed seed and divided into 70\% train, 10\% validation, and 20\% test partitions. The main benchmark uses impaired-arm recordings. The transfer-learning experiments use healthy-arm recordings for pretraining and impaired-arm recordings for finetuning.

\subsection{Model groups}

The experiments are organized into three groups. The first group compares direct LSTM, CNN, and GNN baselines trained under the same preprocessing and feature representation. The second group evaluates optimized CNN students and ResNet teachers, emphasizing the trade-off between multilabel performance and model footprint. The third group evaluates healthy-to-impaired transfer learning. Distillation is implemented in the codebase as a teacher-student compression path, but it is not reported as a completed benchmark in the current results.

\subsection{Evaluation protocol}

The task is five-label multilabel classification with a nominal decision threshold of 0.5 on the held-out test split. We report subset accuracy, also called exact match ratio, mean finger accuracy, macro precision, macro recall, macro F1, macro AUROC, macro AUPRC, and per-finger metrics. Subset accuracy is strict because all five labels must be correct for a sample to count as correct; finger accuracy and macro metrics expose behavior that subset accuracy alone can hide. The present tables report single deterministic train/validation/test split estimates from committed artifacts, so small differences should be interpreted as empirical comparisons rather than statistical significance claims.

The Results section follows the experimental sequence: direct baselines, optimized CNN students and ResNet teachers, transfer learning, finger-level behavior, and literature positioning.
